\chapter{常用微分方程解法}

物理学中的许多定律都可以写成微分方程。本附录总结最常见的几类常微分方程及其标准求解步骤。

\section{可分离变量方程}

标准形式为
\[
\dv{y}{x}=g(x)h(y).
\]

若 $h(y)\neq 0$，可整理为
\[
\frac{1}{h(y)}\dd y=g(x)\dd x,
\]
再两边积分：
\[
\int \frac{1}{h(y)}\dd y=\int g(x)\dd x+C.
\]

\subsection*{典型例子：指数衰减}

放射性衰变、RC 放电、牛顿冷却都常出现
\[
\dv{N}{t}=-\lambda N.
\]

分离变量得
\[
\frac{\dd N}{N}=-\lambda\dd t,
\]
因此
\[
N(t)=N_0 e^{-\lambda t}.
\]

\section{一阶线性 ODE（积分因子法）}

标准形式为
\[
\dv{y}{x}+P(x)y=Q(x).
\]

取积分因子
\[
\mu(x)=e^{\int P(x)\dd x}.
\]

原方程同乘 $\mu(x)$ 后，可化为
\[
\dv{}{x}\qty[\mu(x)y]=\mu(x)Q(x).
\]

积分得到通解：
\[
 y=\frac{1}{\mu(x)}\qty(\int \mu(x)Q(x)\dd x+C).
\]

\subsection*{典型例子：受迫一阶驰豫}

例如电容充电方程
\[
\dv{q}{t}+\frac{1}{RC}q=\frac{V_0}{R}.
\]

积分因子为 $\mu(t)=e^{t/(RC)}$，故
\[
q(t)=CV_0+\qty(q_0-CV_0)e^{-t/(RC)}.
\]

\section{二阶常系数齐次 ODE（特征方程法）}

标准形式为
\[
 a\dv[2]{y}{x}+b\dv{y}{x}+cy=0.
\]

设试探解 $y=e^{rx}$，代入得特征方程
\[
 ar^2+br+c=0.
\]

根据根的类型分三种情况：

\subsection*{（1）两个不同实根 $r_1\neq r_2$}
\[
 y=C_1 e^{r_1 x}+C_2 e^{r_2 x}.
\]

\subsection*{（2）重根 $r_1=r_2=r$}
\[
 y=(C_1+C_2 x)e^{rx}.
\]

\subsection*{（3）共轭复根 $r=\alpha\pm i\beta$}
\[
 y=e^{\alpha x}\qty(C_1\cos\beta x+C_2\sin\beta x).
\]

\subsection*{典型例子：简谐振动}

无阻尼弹簧振子满足
\[
 m\dv[2]{x}{t}+kx=0.
\]

特征方程为 $mr^2+k=0$，故
\[
 r=\pm i\sqrt{\frac{k}{m}}=\pm i\omega,
\qquad
\omega=\sqrt{\frac{k}{m}}.
\]

因此解为
\[
 x(t)=A\cos \omega t+B\sin \omega t.
\]

\begin{figure}[htbp]
\centering
\begin{tikzpicture}
\begin{axis}[
  width=0.65\textwidth, height=0.55\textwidth,
  xlabel={位移 $x$}, ylabel={速度 $v$},
  xmin=-1.2, xmax=1.2, ymin=-1.7, ymax=1.7,
  axis lines=middle,
  grid=major, grid style={gray!30},
  thick,
]
\addplot[blue, domain=0:360, samples=200] ({cos(x)},{-1.5*sin(x)});
\end{axis}
\end{tikzpicture}
\caption{简谐振子在相平面中的轨迹是一条闭合椭圆}
\end{figure}

\subsection*{典型例子：阻尼振动}

阻尼振子满足
\[
\dv[2]{x}{t}+2\beta\dv{x}{t}+\omega_0^2 x=0.
\]
其特征方程为
\[
 r^2+2\beta r+\omega_0^2=0.
\]
因此按阻尼强弱分为三类：
\begin{itemize}
  \item 若 $\beta<\omega_0$，为欠阻尼，解呈现振荡并带有指数衰减包络；
  \item 若 $\beta=\omega_0$，为临界阻尼，系统恰好不振荡且最快回到平衡；
  \item 若 $\beta>\omega_0$，为过阻尼，解是两个指数衰减项的叠加，回到平衡更慢。
\end{itemize}

\begin{figure}[htbp]
\centering
\begin{tikzpicture}
\begin{axis}[
  width=0.75\textwidth, height=0.5\textwidth,
  xlabel={无量纲时间 $\omega_0 t$}, ylabel={位移比 $x/x_0$},
  xmin=0, xmax=6, ymin=-0.3, ymax=1.05,
  axis lines=middle,
  grid=major, grid style={gray!30},
  thick,
  legend style={at={(0.98,0.98)}, anchor=north east, draw=none, fill=none}
]
\addplot[blue, domain=0:6, samples=200] {exp(-0.25*x)*(cos(deg(0.9682458*x)) + 0.2581989*sin(deg(0.9682458*x)))};
\addlegendentry{欠阻尼}
\addplot[red, dashed, domain=0:6, samples=200] {(1 + x)*exp(-x)};
\addlegendentry{临界阻尼}
\addplot[green!60!black, dashdotted, domain=0:6, samples=200] {1.1708204*exp(-0.381966*x) - 0.1708204*exp(-2.618034*x)};
\addlegendentry{过阻尼}
\end{axis}
\end{tikzpicture}
\caption{阻尼振子的三类典型解随时间的衰减形态}
\end{figure}

\section{二阶常系数非齐次 ODE（待定系数法）}

标准形式为
\[
 a\dv[2]{y}{x}+b\dv{y}{x}+cy=f(x).
\]

求解思路：
\begin{enumerate}
  \item 先求对应齐次方程的通解 $y_h$；
  \item 再根据外力项 $f(x)$ 的形式猜测一个特解 $y_p$；
  \item 总解为 $y=y_h+y_p$。
\end{enumerate}

\subsection*{常见试探形式}

\begin{center}
\begin{tabular}{ll}
\toprule
$f(x)$ 的形式 & 可选特解形式 \\
\midrule
多项式 $P_n(x)$ & 同阶多项式 \\
$e^{\alpha x}$ & $Ae^{\alpha x}$ \\
$\sin \beta x$ 或 $\cos \beta x$ & $A\cos \beta x+B\sin \beta x$ \\
$e^{\alpha x}P_n(x)$ & $e^{\alpha x}$ 乘同阶多项式 \\
$e^{\alpha x}\cos \beta x$、$e^{\alpha x}\sin \beta x$ & $e^{\alpha x}(A\cos \beta x+B\sin \beta x)$ \\
\bottomrule
\end{tabular}
\end{center}

若试探特解与齐次解重复，则需乘以 $x$（必要时乘以更高次幂）以保持线性无关。

\subsection*{典型例子：恒力驱动}

方程
\[
\dv[2]{x}{t}+\omega^2 x=\frac{F_0}{m}
\]
对应齐次解为
\[
 x_h=A\cos\omega t+B\sin\omega t.
\]

因为右端是常数，可设特解 $x_p=C$，代入得
\[
 \omega^2 C=\frac{F_0}{m}
 \quad\Rightarrow\quad
 C=\frac{F_0}{m\omega^2}.
\]

故总解为
\[
 x(t)=A\cos\omega t+B\sin\omega t+\frac{F_0}{m\omega^2}.
\]

\section{常见物理微分方程总结表}

\begin{longtable}{p{0.22\textwidth}p{0.32\textwidth}p{0.38\textwidth}}
\toprule
物理情景 & 微分方程 & 通解或关键结论 \\
\midrule
\endfirsthead

\toprule
物理情景 & 微分方程 & 通解或关键结论 \\
\midrule
\endhead

\bottomrule
\endfoot

自由衰减与冷却 & $\dv{y}{t}=-ky$ & $y(t)=y_0 e^{-kt}$ \\

受恒定源驱动的一阶过程 & $\dv{y}{t}+ky=q$ & $y=y_{\infty}+(y_0-y_{\infty})e^{-kt}$，其中 $y_{\infty}=q/k$ \\

简谐振动 & $\dv[2]{x}{t}+\omega^2 x=0$ & $x=A\cos\omega t+B\sin\omega t$ \\

阻尼振动 & $\dv[2]{x}{t}+2\beta\dv{x}{t}+\omega_0^2 x=0$ & 由特征根决定：欠阻尼、临界阻尼、过阻尼三类 \\

受迫振动 & $\dv[2]{x}{t}+2\beta\dv{x}{t}+\omega_0^2 x=f(t)$ & 总解为齐次解加特解；周期驱动下会出现共振 \\

RC 电路充放电 & $R\dv{q}{t}+\frac{1}{C}q=V(t)$ & 若 $V$ 为常数，则电荷以时间常数 $RC$ 指数逼近平衡值 \\

RL 电路电流增长 & $L\dv{I}{t}+RI=\mathcal{E}$ & $I(t)=\frac{\mathcal{E}}{R}+\qty(I_0-\frac{\mathcal{E}}{R})e^{-Rt/L}$ \\

牛顿第二定律的一维形式 & $m\dv[2]{x}{t}=F(x,t,v)$ & 根据受力模型选用分离变量、线性法或数值法 \\

\end{longtable}

\section{解题提醒}

\begin{itemize}
  \item 先辨别方程类别，再选方法，不要一上来就盲目积分。
  \item 求出通解后，务必利用初始条件或边界条件确定积分常数。
  \item 在物理问题中，指数衰减、简谐振动与受迫振动是最常见的三大原型。
\end{itemize}
